
\section{Twelfth Sunday after Trinity: A Man Born Blind}

\begin{quote}

John 9:24–38. Then they called the man who had been blind a second time and said to him: Give God the glory! We know that this man is a sinner. Then he answered: Whether he is a sinner, I do not know; one thing I do know: I was blind, and now I see. So they said to him again: What did he do to you? How did he open your eyes? He answered them: I have told you already, and you did not listen. Why do you want to hear it again? Do you also want to become his disciples? Then they reviled him and said: You are his disciple, but we are disciples of Moses. We know that God has spoken to Moses; but as for this man, we do not know where he is from. The man answered them: "This is indeed a strange thing: that you do not know where he comes from, and yet he has opened my eyes." We know that God does not hear sinners; but if anyone is God-fearing and does his will, him he hears. Never since the world began has it been heard that anyone opened the eyes of one born blind. If this man were not from God, he could do nothing at all. They answered him: You were born wholly in sin, and do you teach us? And they cast him out. Jesus heard that they had cast him out, and when he found him, he said to him: Do you believe in the Son of God? He answered: Who is he, Lord, that I may believe in him? Jesus said to him: You have both seen him; it is he who is speaking with you. Then he said: I believe, Lord. And he worshiped him.

\end{quote}

\bigskip

This passage, like so many others in the New Testament, deals with how a soul found its way to Jesus and was saved. Innumerable are the outward means the Father uses, by his Holy Spirit, to draw people to the Savior, according to the rule the triune God has set for himself, that “no one can come to the Son unless the Father has drawn him” (John 6:44); in this case, as so often, it happened through the miraculous healing of a bodily infirmity, as the man born blind had his eyes opened by Jesus, so that his spiritual blindness might be removed and he might see that Jesus is the Christ.

“What do you want me to do for you?” Jesus asks blind Bartimaeus (Mark 10:51). “Rabboni, that I may receive my sight,” he answered, was healed, and followed Jesus.

That is how it happened there, and that is how it happens here. What matters most is not outward healing, but to be given sight to see Jesus and follow him. So it happened then and so it happens now: one is saved, while another is left behind, and the one whose eyes are opened is saved, but scarcely, and as through fire.

First of all, we note that here only one is spoken of as having been saved; it was the man born blind, who was first healed bodily, but thereafter also became spiritually seeing.

There were many, both neighbors and others, who were witnesses to the miracle that had taken place. Were they saved because of that? Certainly not; no more than if all God’s children in our day were given power to work miracles; scarcely one more soul would be saved.

On the contrary. We see that they were astonished. Then we see them begin to discuss whether it really was a miracle that had taken place—exactly as it would go in our highly enlightened times. Finally, when the man born blind had explained both by whom and how it had happened, we see that they brought the healed man before the Pharisees—almost as though he had committed some evil. For it had happened on a Sabbath. There was a loophole for the conscience! How many thousands of excuses does not the human heart have in order not to see God’s miracles and repent?

Then there are his parents. One would surely think they would from the heart give thanks to God for the unspeakable mercy that their child had become seeing; they ought surely to repent and believe in the Son. But no; they feared men more than God, and since Jesus was mixed up in this matter, and since expulsion from the synagogue had already been decreed for confessing him as Christ, they excused themselves and withdrew. They surely wanted to have nothing to do with so dangerous a man. They wanted to keep themselves safe. A frequent cause of the damnation of thousands upon thousands! They do not have the courage to break with the world and follow Jesus.

Then finally there are those who had the reputation of being the most earnest and God-fearing among the Jews, the Pharisees. They not only did not repent, but they hated Jesus and took offense at the miracle and at the one on whom it had been done.

They had first questioned the healed man born blind once already and examined him closely. Nor were they at once fully agreed. Some thought immediately that a man who did not keep the Sabbath could not be from God, while others, who could not immediately stifle the voice of conscience, wondered whether in general “a sinful man could do such signs.” But once it was beyond all doubt that the man really had been blind from birth and really had been healed, so that here too there was no loophole for not acknowledging Jesus, by whom the miracle had been done, then they were soon fully agreed.

It was then that they called the man born blind a second time and said: “Give God the glory! We know that this man is a sinner.”

From this we see that one may be both learned and well-versed in God’s Word, and pious in speech, without having a true and living faith in Jesus.

The miracle, which they could not deny, had made them so zealous for God and his glory that with so much the better conscience they could denounce the one through whom the miracle had been done, the Sabbath-breaker. Their great words, “we know,” were meant to put an end to all doubt and all objections in the simple man born blind, and thereby be rid of this whole unpleasant matter.

But they did not know that the one whose eyes Jesus opens sees what is often hidden from the wise and intelligent. He could not follow such logic: to honor God and blame the one through whom he works miracles. Instead of any other answer, he simply points to what has happened: “One thing I know, that I, who was blind, now see.”

Oh, what a glorious, mighty, and irresistible confession! It is like when the simple Christian can bear witness: “One thing I know, that I am born again and have the witness of God’s Spirit in my spirit, that I am a child of God!”

No wonder they were thrown into confusion and began again to ask how the whole thing had happened, and that the simple man born blind could not refrain from taunting these most holy and learned Pharisees with their own groundless hatred of Jesus and asking them whether perhaps they too wished to become his disciples.

The high leaders of the people answered by wrapping themselves in their immense dignity: We are Moses’ disciples. To Moses God has spoken, but this man we do not know where he is from. Oh, what proud blindness!

Jesus had rightly said of them: “If you believed Moses, you would believe me also; for he wrote of me,” as we also sing:

“They search the Scripture closely, But Jesus is not known, Though he stands before their eyes, For Pride says: No.”

That the man born blind had not only had his bodily eyes opened, but had also received a new and simple sight into God’s Word, his answer brings it to light.

A divine miracle has been done upon me, and you yourselves have said: Give God the glory! How then is it possible that you, teachers of Israel, do not know where he is from, and indeed even call him a sinner? God, after all, does not hear sinners.

Already the man born blind is driven to contend for his benefactor; it is given him what he shall say, and he is ready to take the consequences, in other words to endure scandal and reproach, two strong witnesses that a man is truly on the way of repentance.

The Pharisees feel the truth and the sting of this answer, and they seized the weapon which in that time and in all times has belonged to them. They reviled him and cast him out.

How many simple children of God down here on earth had to suffer the same treatment?

But now the sufferings began in earnest. Cast out of the synagogue. That means: forsaken, rejected by all. No wonder the man born blind now groans with a deeper cry: Is there no physician in Gilead? Where is the Son of God? Where is the Messiah?

Poor anxious and troubled soul! Though people have forgotten you and cast you off, you are not forgotten by him who goes out into the wilderness to seek his lost sheep.

When Jesus heard that the man born blind had been cast out, then he came to him. That is what is written.

“Do you believe in the Son of God?” he asked. He knew how far the man born blind had come; he had seen to what point he had come: that he was earnestly seeking grace.

And so he answers, as if he meant to say: Is he not the very one I seek, is he not the one I have been waiting for? Where is he? Where is he?

Friend, have you thus sought and thus asked and—received this answer:

Here he is! Here he is! Right beside you—in the Word.

Then the man born blind had his eyes opened fully; then he received the anointing of the Holy Spirit and knew all things; for he knew Jesus, yes, much more, he was known by him, fell down before him and worshiped him: I believe, Lord!

Oh, what a blessed moment! To be born again, to believe in Jesus, to call him brother and God Father!

Wonder upon wonder, that I, a natural man, who do not grasp the things that belong to the Spirit of God, should have my eyes opened by faith in the Son of God and see God and not die, but live!